
\section{Fundamentos de modelado y simulación}
\label{subsec:fundamentos_modelado_simulacion}

Para el estudio de sistemas cuya observación directa es insuficiente, arriesgada, limitada o costosa para examinar todas las condiciones de interés y relevantes, el modelado y la simulación se convierten en metodologías fundamentales. Con estas metodologías se crean representaciones simplificadas de un acontecimiento, se determinan conexiones entre sus elementos y se realizan experimentos computacionales bajos condiciones controladas. Por lo tanto, la simulación no se limita a la representación visual de un entorno, sino que permite formular preguntas, modificar variables, comparar configuraciones y analizar las consecuencias de diferentes decisiones sin ejercer una intervención directa sobre el sistema real \cite{bungartz2014modeling,law2015simulation}. 

\subsection{Sistema, modelo y simulación}
\label{subsubsec:sistema_modelo_simulacion}

\subsubsection{Sistema}

Se define un sistema como un grupo de elementos interrelacionados entre sí y, de manera conjunta, generan un comportamiento observable. Para analizar un sistema, es fundamental definir sus límites, identificar los componentes que forman parte directa de su funcionamiento, reconocer las variables de entrada que modifican su estado y determinar las salidas o resultados que se quieren observar \cite{bungartz2014modeling, law2015simulation}. 

La delimitación del sistema resulta fundamental, ya que un mismo fenómeno puede analizarse desde distintas perspectivas. En el caso del tráfico urbano, por ejemplo, los vehículos, la infraestructura vial, los semáforos, las zonas de origen y destino y las reglas de circulación pertenecen a los componentes internos del sistema, mientras que las variables como la demanda de viajes, o la duración de la corrida pueden tratarse como condiciones externas o parámetros del experimento.

\subsubsection{Modelo}

Un modelo es una representación deliberadamente simplificada de un sistema, construida con el propósito de describir, explicar, predecir o experimentar con determinados aspectos de su comportamiento. En este sentido, un modelo no incorpora necesariamente todos los elementos del sistema real, sino aquellos que son pertinentes para responder las preguntas del estudio \cite{law2015simulation}. 

La simplificación constituye una condición necesaria del modelado, debido a que una representación que intenta reproducir cada detalle del fenómeno perdería utilidad práctica y aumentaría de forma considerable su complejidad computacional. Por otra parte, una simplificación excesiva podría eliminar relaciones relevantes y producir resultados poco representativos. Por esta razón, el nivel de detalle debe establecerse de acuerdo con el objetivo del modelo, las variables que se pretenden analizar y las limitaciones de información y capacidad computacional \cite{law2015simulation}.

El proceso de modelado transforma una situación real o concepto abstracto en una estructura que puede analizarse a través de reglas, relaciones matemáticas, algoritmos o procedimientos computacionales. Al menos tres niveles pueden diferenciarse a partir de eso: el modelo computacional, el modelo conceptual y el sistema de referencia. El sistema de referencia corresponde al fenómeno, proceso o realidad que se desea estudiar y sobre el cual se delimitan los componentes, condiciones y comportamientos de interés. Por su parte, el modelo conceptual describe de manera general los componentes, supuestos, relaciones y límites del fenómeno, mientras que el modelo computacional traduce dicha descripción en estructuras de datos, parámetros y algoritmos ejecutables \cite{bungartz2014modeling,law2015simulation}. 

\subsubsection{Simulación}

La simulación corresponde a la ejecución de un modelo para observar su comportamiento bajo determinadas condiciones. Por tanto, modelo y simulación no son términos equivalentes: el modelo establece qué elementos y relaciones se representan, mientras que la simulación aplica esas reglas durante un intervalo, una secuencia de eventos o un conjunto de escenarios \cite{banks2010discrete,law2015simulation}. Una corrida de simulación representa una ejecución particular determinada por su configuración inicial, sus parámetros, sus condiciones de operación y, cuando existe aleatoriedad, la semilla utilizada. En consecuencia, varias corridas pueden emplear el mismo modelo y producir resultados diferentes si se modifican las condiciones iniciales o las variables aleatorias \cite{law2015simulation}.

\subsection{Clasificación de modelos}
\label{subsubsec:clasificacion_modelos}

Los modelos pueden clasificarse de acuerdo con la forma en que representan el tiempo, la incertidumbre, los cambios de estado y el nivel de detalle del sistema. Estas clasificaciones no son mutuamente excluyentes; por el contrario un mismo modelo puede pertenecer simultáneamente a varias categorías \cite{bungartz2014modeling,law2015simulation}. Por lo tanto, identificar la naturaleza del modelo permite establecer qué técnicas de implementación, experimentación y análisis son apropiadas, así como reconocer las limitaciones derivadas de la representación seleccionada. 

Según su relación con el tiempo, los modelos se dividen en estáticos y dinámicos. Un modelo estático es aquel que representa el sistema en un instante o mediante relaciones que no describen explícitamente su evolución temporal. Por otra parte, un modelo dinámico considera que el estado del sistema cambia a medida que transcurre el tiempo \cite{law2015simulation}. 

De acuerdo con el tratamiento de la incertidumbre, los modelos pueden ser determinísticos o estocásticos. En un modelo determinístico, una misma configuración inicial produce las mismas salidas siempre que no se alteren sus parámetros ni sus reglas. Por otro lado, un modelo estocástico incorpora variables aleatorias o distribuciones de probabilidad, por lo que una misma configuración general puede producir distintas realizaciones \cite{bungartz2014modeling,law2015simulation}. 

Asimismo, los modelos se clasifican como continuos, discretos o híbridos. En los modelos continuos, las variables de estado pueden cambiar de manera continua respecto del tiempo. En los modelos discretos, los cambios se representan en instantes determinados o como consecuencia de eventos específicos. Por último, los modelos híbridos combinan ambos enfoques \cite{banks2010discrete,bungartz2014modeling}. 

\begin{comment}
    Esto podria agregarlo hasta después que estoy mencionando CARLA
CARLA representa movimiento, velocidad y dinámica vehicular mediante modelos físicos, pero su ejecución computacional avanza a través de pasos de simulación o \textit{ticks}. Por consiguiente, el experimento se observa y registra en intervalos discretos, aunque algunas variables representen fenómenos físicamente continuos.
\end{comment}

Finalmente, los modelos de simulación de tráfico pueden diferenciarse según su nivel de agregación en macroscópicos, mesoscópicos y microscópicos. Los modelos macroscópicos representan el tráfico mediante variables agregadas, como el flujo, la densidad y la velocidad media, sin describir individualmente cada vehículo. Por otra parte, los modelos microscópicos representan los vehículos de manera individual y consideran sus interacciones con otros actores y con la infraestructura vial. Los modelos mesoscópicos adoptan un nivel intermedio, debido a que pueden conservar la representación de vehículos individuales o grupos de vehículos, pero simplifican sus interacciones y reglas de comportamiento \cite{burghout2004hybrid}. 

% Hasta acá todo nítido de momento

\subsection{Verificación, validación y reproducibilidad}
\label{subsubsec:verificacion_validacion_reproducibilidad}

La calidad de una simulación no depende únicamente de que el programa se ejecute sin interrupciones. También es necesario determinar si el modelo conceptual resulta adecuado para el propósito del estudio, si la implementación representa correctamente ese modelo y si los resultados poseen suficiente consistencia para sustentar el análisis. En este contexto, la verificación y la validación constituyen procesos relacionados, pero conceptualmente diferentes \cite{sargent2003verification}. 

La verificación busca determinar si el modelo computacional implementa correctamente el modelo conceptual. En términos prácticos, responde a la pregunta de si el sistema se construye de acuerdo con sus especificaciones \cite{sargent2003verification}. Este proceso puede incluir revisión de código, pruebas unitarias, inspección de archivos de configuración, comprobación de rangos, validación de rutas, seguimiento de estados y análisis de registros. 



\begin{comment}
    En el generador de tráfico, la verificación comprende comprobar que los vehículos se crean en puntos válidos, que las probabilidades suman la unidad, que los destinos se seleccionan según la distribución configurada, que las rutas son transitables y que los datos registrados corresponden con el estado observado en la simulación.
\end{comment}

La validación, por otra parte, busca establecer si el modelo representa con suficiente fidelidad el sistema de interés para el uso previsto. Sargent distingue aspectos como la validez del modelo conceptual, la validez de los datos y la validez operacional de los resultados \cite{sargent2003verification}. En consecuencia, la validación no debe interpretarse como una demostración de que el modelo reproduce toda la realidad, sino como una evaluación de su adecuación para responder preguntas delimitadas. 

La validación, por otra parte, busca establecer si el modelo representa con suficiente precisión el sistema de interés para el uso previsto. Sargent distingue entre la validez del modelo conceptual, la validez de los datos y la validez operacional, además de separar estos procesos de la verificación del modelo computacional \cite{sargent2003verification}. Por esta razón, la validación no debe interpretarse como una demostración de que el modelo reproduce todos los aspectos de la realidad, sino como una evaluación de su adecuación para responder preguntas delimitadas dentro de un dominio de aplicación específico.

\begin{comment}
En este proyecto, la validez se analiza dentro del dominio operacional definido (ODD, por sus siglas en inglés) y se relaciona con la capacidad del generador para producir viajes estructurados, rutas ejecutables, distribuciones coherentes con las probabilidades establecidas y registros útiles para el análisis preliminar de movilidad.
\end{comment}

La reproducibilidad se refiere a la posibilidad de obtener resultados computacionales consistentes cuando se utilizan los mismos datos de entrada, código, métodos, pasos de procesamiento y condiciones de análisis. Por otra parte, la replicabilidad implica desarrollar un nuevo estudio dirigido a responder la misma pregunta mediante datos obtenidos de forma independiente \cite{nasem2019reproducibility}. 

\begin{comment}
Para favorecer la reproducibilidad, cada corrida debe conservar información suficiente sobre la versión de CARLA, el mapa utilizado, la semilla, el número de vehículos, la duración, el paso temporal, las probabilidades, las zonas habilitadas, las reglas del Traffic Manager y los archivos de entrada. Además, el código y los datos generados deben asociarse con un identificador de ejecución y almacenarse de manera trazable. Las National Academies señalan que la reproducibilidad computacional requiere descripciones claras de los métodos, el código, los datos y el flujo de análisis \cite{nasem2019reproducibility}. Por consiguiente, el registro de estos elementos no constituye únicamente una decisión de implementación, sino una condición metodológica para comparar corridas y sustentar los resultados.
\end{comment}

\begin{comment}
\subsubsection{Integración conceptual del marco teórico}
\label{subsubsec:integracion_conceptual_marco}

El diagrama conceptual de movilidad vehicular, modelos de CARLA y evaluación del generador se incorpora al finalizar los fundamentos de modelado y simulación porque sintetiza la transición entre el fenómeno de interés y su representación computacional. Su estructura parte de la movilidad vehicular como tema general, continúa con los sistemas de transporte y las interacciones del tráfico urbano, incorpora los modelos de movilidad y elección de destinos, y conduce hacia la simulación mediante CARLA. A partir de ello, el diagrama relaciona el objetivo del generador con el análisis de las métricas y con la representación de los resultados. Por tanto, su propósito no es describir la arquitectura de implementación, sino establecer la lógica conceptual que organiza el marco teórico.

\begin{figure}[htbp]
    \centering
    \includegraphics[width=\textwidth]{figuras/diagrama_conceptual_movilidad_carla.pdf}
    \caption{Diagrama conceptual de movilidad vehicular, modelos de CARLA y evaluación del generador. Fuente: elaboración propia.}
    \label{fig:diagrama_conceptual_movilidad_carla}
\end{figure}

Como se muestra en la Figura \ref{fig:diagrama_conceptual_movilidad_carla}, el tema general corresponde a la movilidad vehicular, mientras que el sistema se define por la interacción entre vehículos, infraestructura y reglas de circulación. Los modelos proporcionan las abstracciones necesarias para representar viajes, tráfico, elección de destinos y rutas; posteriormente, CARLA permite ejecutar estas abstracciones dentro de un entorno simulado. El objetivo se concentra en desarrollar un generador de tráfico urbano verosímil y medible dentro de un ODD, mientras que el análisis utiliza la distribución de destinos, el tiempo de viaje, la velocidad media y las incidencias para comparar escenarios. Finalmente, la representación mediante tablas y gráficas permite interpretar los resultados sin atribuir al proyecto funciones de percepción, decisión autónoma integral o control de un vehículo real.

\end{comment}

\begin{comment}
\subsubsection*{Fuentes utilizadas por párrafo}

\textbf{Párrafo 18. Propósito del diagrama conceptual}

Elaboración propia. El párrafo organiza los conceptos del marco teórico según el alcance y la secuencia lógica del proyecto. No introduce una definición tomada de una fuente externa.

\textbf{Párrafo 19. Interpretación del diagrama conceptual}

Elaboración propia. La relación entre tema, sistema, modelo, simulación, objetivo, análisis y representación sintetiza la estructura conceptual propuesta para el trabajo.
\end{comment}


\begingroup
\let\seccionoriginal\section
\let\subseccionoriginal\subsection
\renewcommand{\section}[1]{\seccionoriginal{#1}\mbox{}\par}
\renewcommand{\subsection}[1]{\subseccionoriginal{#1}\mbox{}\par}

\section{Modelos basados en agentes}
\label{subsec:modelos_basados_agentes}

Los modelos basados en agentes representan un sistema a partir de entidades individuales que poseen estados, reglas de comportamiento y mecanismos de interacción. A diferencia de los enfoques agregados, que describen el comportamiento global mediante variables resumidas, este tipo de modelado permite observar cómo las acciones locales de múltiples entidades producen patrones colectivos. Por consiguiente, resulta adecuado para estudiar fenómenos donde la interacción entre unidades individuales influye en el comportamiento general del sistema, como ocurre en la movilidad vehicular urbana \cite{russell1995ai,macal2010abms,bonabeau2002abm}.

\subsection{Definición de agente}
\label{subsubsec:definicion_agente}

Un agente puede definirse como una entidad capaz de percibir un entorno, procesar información relevante y ejecutar acciones que modifican su propio estado o el estado del sistema en el que participa. Desde la perspectiva de la inteligencia artificial, esta relación se describe mediante percepciones, acciones, objetivos y entorno; es decir, el agente recibe información, selecciona una acción y actúa bajo una medida de desempeño o un criterio de éxito \cite{russell1995ai}. Esta definición no exige que el agente sea completamente inteligente ni que posea aprendizaje avanzado, sino que tenga una lógica de comportamiento identificable dentro del modelo.

En simulación, un agente no necesariamente corresponde a una entidad física real. También puede ser una representación computacional de un vehículo, peatón, semáforo, usuario o componente operativo del sistema. Lo relevante es que dicha entidad cuente con atributos propios, reglas de decisión y capacidad de interacción con otros elementos del entorno. De esta manera, los modelos basados en agentes permiten representar heterogeneidad, ya que cada agente puede diferenciarse por su posición, velocidad, destino, ruta, estado operativo o comportamiento asignado \cite{macal2010abms,wooldridge2009multiagent}.

\subsection{Agentes vehiculares}
\label{subsubsec:agentes_vehiculares}

Un agente vehicular representa a un vehículo individual dentro de una red vial simulada. Este agente posee un estado dinámico compuesto por variables como posición, orientación, velocidad, carril, ruta asignada, destino, tiempo de salida y estado del viaje. Además, puede seguir reglas de circulación, mantener distancia con otros vehículos, reaccionar ante señales de tránsito, ajustar su velocidad y avanzar de acuerdo con una trayectoria permitida por la infraestructura vial. Por otra parte, el agente vehicular también puede vincularse con decisiones de mayor nivel, como la selección de un destino o la asignación de una ruta.

En \textit{Car Learning to Act} (CARLA), los vehículos se representan como actores dentro del mundo simulado y pueden ser controlados mediante la API de Python o mediante el \textit{Traffic Manager}. Este último permite controlar vehículos en modo \textit{autopilot}, configurar ciertos parámetros de comportamiento y poblar la simulación con tráfico urbano configurable \cite{dosovitskiy2017carla,carla_traffic_manager_docs}. Para el generador propuesto, esta distinción es importante: la lógica estocástica define qué destino se selecciona y qué ruta se asigna, mientras que CARLA y el \textit{Traffic Manager} ejecutan el desplazamiento local del vehículo dentro del escenario.

El agente vehicular utilizado en este proyecto no debe interpretarse como un sistema completo de conducción autónoma. Su finalidad no es integrar percepción avanzada, aprendizaje profundo, toma de decisiones ética ni control físico de un vehículo real. Más bien, se utiliza como una unidad simulada de movilidad que permite estudiar viajes dentro de un escenario urbano controlado. A partir de ello, el vehículo funciona como agente porque posee un origen, un destino, una ruta, un estado de movimiento y una interacción con otros actores de la simulación.

\subsection{Interacción local y comportamiento emergente}
\label{subsubsec:interaccion_local_comportamiento_emergente}

La interacción local ocurre cuando el comportamiento de un agente depende de las condiciones próximas del entorno. En tráfico vehicular, estas condiciones incluyen la presencia de vehículos delanteros, la disponibilidad de carriles, la existencia de semáforos, la geometría de la vía, los límites de velocidad y los posibles conflictos en intersecciones. Aunque cada agente toma decisiones limitadas a su situación inmediata, la suma de muchas interacciones puede producir efectos agregados como colas, disminución de velocidad, concentración de viajes, bloqueos temporales o variación en los tiempos de recorrido \cite{bonabeau2002abm,macal2010abms}.

Este fenómeno se conoce como comportamiento emergente, porque el patrón global no se programa necesariamente de forma directa, sino que aparece como resultado de múltiples interacciones individuales. Por ejemplo, una alta concentración de vehículos hacia una misma zona puede incrementar los tiempos de viaje aunque ningún vehículo tenga una regla explícita que indique "generar congestión". Del mismo modo, pequeños cambios en la distribución de destinos pueden modificar la ocupación de ciertas rutas, la velocidad media y la probabilidad de incidentes dentro del escenario simulado. Por lo que, los modelos basados en agentes permiten conectar decisiones individuales con métricas agregadas de movilidad.

\subsection{Limitaciones}
\label{subsubsec:limitaciones_modelos_basados_agentes}

Los modelos basados en agentes dependen fuertemente de las reglas, parámetros y supuestos utilizados para representar el comportamiento individual. Si las reglas son demasiado simples, el modelo puede omitir aspectos relevantes del comportamiento humano o de la circulación urbana. Si las reglas son excesivamente complejas, el modelo puede volverse difícil de verificar, calibrar e interpretar. Además, los patrones emergentes pueden ser sensibles a las condiciones iniciales, a la cantidad de agentes, a las semillas pseudoaleatorias y a la configuración del entorno \cite{macal2010abms,bonabeau2002abm}.

Otra limitación importante es que la verosimilitud del comportamiento simulado no equivale automáticamente a validez externa. Un agente vehicular puede circular de forma plausible dentro de CARLA y, aun así, no representar con exactitud el comportamiento de un conductor real. Esto se debe a que la simulación abstrae múltiples factores, como estilos de conducción, percepción humana, toma de decisiones bajo presión, errores, fatiga, distracción y variaciones culturales en la conducción. Por esta razón, los resultados deben interpretarse dentro del dominio operacional definido (ODD, por sus siglas en inglés) y bajo las condiciones experimentales configuradas.

Además, CARLA y el \textit{Traffic Manager} facilitan la creación de tráfico urbano controlado, pero no sustituyen por sí mismos un modelo completo de demanda de viajes. La plataforma permite ejecutar vehículos, controlar ciertos comportamientos y registrar datos, pero la lógica de elección de destinos, categorías, probabilidades y comparación con un escenario base debe definirse mediante el generador implementado. De manera complementaria, las métricas obtenidas describen el comportamiento del escenario simulado y no deben presentarse como predicciones directas de una red vial real sin una calibración posterior con datos empíricos.

\section{Variables microscópicas del tráfico}
\label{subsec:variables_microscopicas_trafico}

El análisis microscópico del tráfico estudia el comportamiento de los vehículos de forma individual, en lugar de representar el flujo vehicular únicamente mediante valores agregados. Bajo este enfoque, cada vehículo posee una posición, una velocidad, una aceleración, una separación respecto de otros vehículos y una trayectoria propia a través del tiempo \cite{immers2005trafficflow,fhwa2001trafficflow}. De manera complementaria, estas variables permiten relacionar el movimiento individual de los agentes vehiculares con métricas posteriores como distancia recorrida, tiempo de viaje, velocidad media, vehículos detenidos e interacciones entre vehículos.

A diferencia del análisis macroscópico, que utiliza variables como flujo, densidad y velocidad media para describir el comportamiento agregado de una corriente vehicular, el análisis microscópico conserva la identidad de cada vehículo. Esta característica resulta fundamental para una simulación basada en agentes, ya que permite registrar el estado de cada vehículo durante una corrida y reconstruir posteriormente su comportamiento. Por consiguiente, las variables microscópicas constituyen la base para interpretar cómo se desplaza cada agente vehicular dentro de una red vial simulada y cómo sus movimientos contribuyen a los patrones globales observados \cite{fhwa2001trafficflow}.

\subsection{Posición y trayectoria}
\label{subsubsec:posicion_trayectoria}

La posición describe la ubicación de un vehículo sobre una vía en un instante determinado. En teoría del flujo vehicular, esta posición suele representarse sobre un eje longitudinal \(x\), orientado en la dirección de circulación. Si el vehículo \(i\) se desplaza a lo largo de una vía, su posición se expresa como una función del tiempo:

\begin{equation}
x_i = x_i(t)
\label{eq:posicion_trayectoria}
\end{equation}

donde \(x_i(t)\) representa la posición longitudinal del vehículo \(i\) en el instante \(t\). La variable \(x_i\) se expresa en metros \((m)\) cuando se trabaja en el Sistema Internacional de Unidades, mientras que \(t\) se expresa en segundos \((s)\). Esta formulación asume que el movimiento se proyecta sobre una dirección principal de avance, por lo que resulta adecuada para analizar el desplazamiento de un vehículo sobre un carril, tramo o ruta. En una simulación urbana más compleja, la posición también puede representarse mediante coordenadas bidimensionales o tridimensionales; sin embargo, la proyección longitudinal permite estudiar con claridad el avance del vehículo sobre su trayectoria.

La trayectoria corresponde a la evolución de la posición de un vehículo a través del tiempo. En un diagrama espacio--tiempo, el eje horizontal representa el tiempo y el eje vertical representa la posición; cada vehículo se muestra como una curva que indica cómo cambia su ubicación durante la simulación \cite{immers2005trafficflow,fhwa2001trafficflow}. Una trayectoria con pendiente elevada representa un vehículo que avanza con mayor velocidad, mientras que una trayectoria casi horizontal indica que el vehículo se desplaza lentamente o permanece detenido. Además, cuando dos vehículos circulan en el mismo carril, sus trayectorias no deben cruzarse, ya que ello implicaría que ambos ocupan la misma posición longitudinal en el mismo instante.

En CARLA, la posición de un actor puede consultarse mediante la API de Python a través del método \texttt{get\_location()}, el cual devuelve la localización del actor en metros. Para efectos de análisis de movilidad, esta localización puede registrarse por cada \textit{tick} de simulación y convertirse en una trayectoria discreta. De esta forma, la trayectoria continua \(x_i(t)\) se aproxima mediante una secuencia de observaciones:

\begin{equation}
\left\{x_i(t_0), x_i(t_1), x_i(t_2), \ldots, x_i(t_n)\right\}
\label{eq:trayectoria_discreta}
\end{equation}

donde \(t_0, t_1, t_2, \ldots, t_n\) representan los instantes de muestreo de la simulación, y \(x_i(t_j)\) representa la posición del vehículo \(i\) en el instante \(t_j\). Esta formulación es una adaptación discreta de la trayectoria continua, adecuada para simulaciones computacionales que avanzan mediante pasos de tiempo. Su utilidad radica en que permite calcular distancia recorrida, velocidad aproximada, detenciones y estado de avance del viaje sin depender únicamente de una observación final.

\subsection{Velocidad y aceleración}
\label{subsubsec:velocidad_aceleracion}

La velocidad describe la variación de la posición de un vehículo respecto del tiempo. En una formulación continua, la velocidad longitudinal del vehículo \(i\) se define como la derivada de su trayectoria:

\begin{equation}
v_i(t)=\frac{dx_i(t)}{dt}
\label{eq:velocidad_continua}
\end{equation}

donde \(v_i(t)\) representa la velocidad del vehículo \(i\) en el instante \(t\), expresada en metros por segundo \((m/s)\), \(x_i(t)\) representa su posición longitudinal en metros \((m)\) y \(t\) representa el tiempo en segundos \((s)\). Esta ecuación indica que la velocidad mide qué tan rápido cambia la posición del vehículo. En el análisis de tráfico, la velocidad permite identificar condiciones de circulación fluida, reducción de velocidad, detenciones y posibles efectos de congestión local \cite{immers2005trafficflow,fhwa2001trafficflow}.

En una simulación computacional, las mediciones se obtienen en intervalos discretos. Por ello, la velocidad también puede aproximarse mediante diferencias finitas:

\begin{equation}
v_i(t_j) \approx \frac{x_i(t_j)-x_i(t_{j-1})}{\Delta t}
\label{eq:velocidad_discreta}
\end{equation}

donde \(v_i(t_j)\) representa la velocidad aproximada del vehículo \(i\) en el instante \(t_j\), \(x_i(t_j)\) y \(x_i(t_{j-1})\) son las posiciones registradas en dos instantes consecutivos, y \(\Delta t=t_j-t_{j-1}\) es el intervalo de muestreo en segundos. Esta expresión se obtiene como una aproximación discreta de la derivada definida en la ecuación \eqref{eq:velocidad_continua}. El supuesto principal es que el intervalo \(\Delta t\) es suficientemente pequeño para que el cambio de posición represente adecuadamente el desplazamiento del vehículo entre dos observaciones consecutivas.

La aceleración describe la variación de la velocidad respecto del tiempo. En forma continua, se expresa como:

\begin{equation}
a_i(t)=\frac{dv_i(t)}{dt}=\frac{d^2x_i(t)}{dt^2}
\label{eq:aceleracion_continua}
\end{equation}

donde \(a_i(t)\) representa la aceleración del vehículo \(i\) en el instante \(t\), expresada en metros por segundo cuadrado \((m/s^2)\). La aceleración positiva indica aumento de velocidad, mientras que la aceleración negativa representa desaceleración o frenado. En tráfico vehicular, esta variable permite identificar maniobras de arranque, frenado, aproximación a intersecciones, respuesta ante un vehículo delantero y variaciones de comportamiento inducidas por la circulación local \cite{immers2005trafficflow}.

De forma discreta, la aceleración puede aproximarse mediante la diferencia entre velocidades consecutivas:

\begin{equation}
a_i(t_j) \approx \frac{v_i(t_j)-v_i(t_{j-1})}{\Delta t}
\label{eq:aceleracion_discreta}
\end{equation}

donde \(a_i(t_j)\) representa la aceleración aproximada del vehículo \(i\) en el instante \(t_j\), \(v_i(t_j)\) y \(v_i(t_{j-1})\) son velocidades registradas o calculadas en instantes consecutivos, y \(\Delta t\) es el paso temporal de medición. Esta ecuación es una adaptación por diferencias finitas de la definición continua de aceleración. En simulación, su precisión depende de la estabilidad del paso temporal y de la calidad de las mediciones de posición o velocidad.

CARLA permite consultar directamente el vector de velocidad de un actor mediante \texttt{get\_velocity()} y el vector de aceleración mediante \texttt{get\_acceleration()}, ambos asociados al último \textit{tick} recibido por el cliente \cite{carla_python_api_0915}. Cuando se requiere obtener la rapidez escalar del vehículo a partir de sus componentes vectoriales, puede utilizarse la norma euclidiana:

\begin{equation}
\left\| \vec{v}_i(t) \right\| =
\sqrt{v_{x,i}(t)^2+v_{y,i}(t)^2+v_{z,i}(t)^2}
\label{eq:norma_velocidad}
\end{equation}

donde \(\left\| \vec{v}_i(t) \right\|\) representa la magnitud de la velocidad del vehículo \(i\), expresada en metros por segundo \((m/s)\), y \(v_{x,i}(t)\), \(v_{y,i}(t)\) y \(v_{z,i}(t)\) son las componentes del vector de velocidad en los ejes \(x\), \(y\) y \(z\), respectivamente. Esta expresión se adapta al formato vectorial utilizado por CARLA y permite convertir la velocidad tridimensional del actor en una magnitud escalar útil para calcular velocidad media, detectar vehículos detenidos y comparar corridas.

\subsection{Separación y \textit{headway}}
\label{subsubsec:separacion_headway}

La separación entre vehículos es una variable microscópica central porque permite describir la distancia disponible entre un vehículo seguidor y el vehículo que circula delante de él. Si el vehículo \(i\) sigue al vehículo \(i-1\) sobre el mismo carril y ambos se desplazan en la dirección creciente del eje \(x\), el \textit{space headway} puede expresarse como:

\begin{equation}
s_i(t)=x_{i-1}(t)-x_i(t)
\label{eq:space_headway}
\end{equation}

donde \(s_i(t)\) representa el \textit{space headway} del vehículo \(i\), expresado en metros \((m)\); \(x_{i-1}(t)\) es la posición longitudinal del vehículo líder; y \(x_i(t)\) es la posición longitudinal del vehículo seguidor. Esta formulación asume que las posiciones se miden respecto de un mismo punto de referencia del vehículo, por ejemplo el parachoques trasero, como se plantea en el tratamiento clásico de trayectorias microscópicas \cite{immers2005trafficflow}. El valor de \(s_i(t)\) incluye la longitud física del vehículo líder y la distancia libre entre ambos vehículos.

La distancia libre o \textit{gap} espacial excluye la longitud del vehículo líder. Por ello, puede definirse como:

\begin{equation}
d_i(t)=s_i(t)-L_{i-1}
\label{eq:gap_espacial}
\end{equation}

donde \(d_i(t)\) representa la distancia libre entre el vehículo seguidor \(i\) y el vehículo líder \(i-1\), expresada en metros \((m)\), y \(L_{i-1}\) representa la longitud física del vehículo líder en metros. Esta variable resulta útil para analizar condiciones de seguimiento, seguridad operacional y proximidad entre vehículos. En un escenario simulado, valores muy bajos de \(d_i(t)\) pueden indicar circulación densa, riesgo de interacción cercana o necesidad de revisar los parámetros de seguimiento del controlador vehicular.

El \textit{time headway} describe la separación temporal entre dos vehículos consecutivos al pasar por un mismo punto de referencia. Si el vehículo líder \(i-1\) pasa por una sección de medición en el instante \(t_{i-1}\) y el vehículo seguidor \(i\) pasa por la misma sección en el instante \(t_i\), entonces:

\begin{equation}
h_i=t_i-t_{i-1}
\label{eq:time_headway}
\end{equation}

donde \(h_i\) representa el \textit{time headway}, expresado en segundos \((s)\). Esta variable se mide en una sección fija de la vía y permite analizar la frecuencia temporal con la que los vehículos atraviesan un punto. En teoría de flujo de tráfico, el \textit{time headway} se relaciona con el flujo vehicular, ya que intervalos temporales pequeños implican que más vehículos pasan por una sección durante un período determinado \cite{fhwa2001trafficflow}.

Bajo el supuesto simplificado de velocidad aproximadamente constante durante un intervalo corto, el \textit{time headway} también puede aproximarse a partir del \textit{space headway}:

\begin{equation}
h_i(t) \approx \frac{s_i(t)}{v_i(t)}
\label{eq:headway_aproximado}
\end{equation}

donde \(h_i(t)\) se expresa en segundos \((s)\), \(s_i(t)\) se expresa en metros \((m)\) y \(v_i(t)\) en metros por segundo \((m/s)\). Esta aproximación supone que el vehículo seguidor mantiene una velocidad casi constante durante el intervalo considerado y que el espaciamiento observado puede traducirse en separación temporal. En tráfico urbano, este supuesto debe utilizarse con cautela porque las aceleraciones, frenadas, intersecciones y semáforos pueden modificar rápidamente la relación entre distancia y tiempo.

\subsection{Interacción vehículo-seguidor}
\label{subsubsec:interaccion_vehiculo_seguidor}

La interacción vehículo-seguidor describe cómo el movimiento de un vehículo depende del vehículo que circula delante de él. Esta relación es uno de los fundamentos de los modelos de seguimiento vehicular, aunque su desarrollo completo corresponde a la sección de modelos microscópicos de comportamiento. Desde el punto de vista de las variables, la interacción puede observarse mediante la separación, la velocidad del seguidor, la velocidad del líder y la diferencia relativa de velocidades \cite{immers2005trafficflow,fhwa2001trafficflow}.

La diferencia de velocidad entre el vehículo seguidor y el vehículo líder se define como:

\begin{equation}
\Delta v_i(t)=v_i(t)-v_{i-1}(t)
\label{eq:diferencia_velocidad}
\end{equation}

donde \(\Delta v_i(t)\) representa la velocidad relativa del vehículo \(i\) respecto del vehículo \(i-1\), expresada en metros por segundo \((m/s)\). Si \(\Delta v_i(t)>0\), el vehículo seguidor se aproxima al vehículo líder porque circula a mayor velocidad. Si \(\Delta v_i(t)<0\), el vehículo seguidor se aleja o aumenta su separación relativa. Esta variable permite interpretar situaciones de aproximación, seguimiento estable o desacoplamiento entre vehículos consecutivos.

La dinámica básica de seguimiento puede describirse mediante una función general:

\begin{equation}
a_i(t)=f\left(v_i(t), \Delta v_i(t), d_i(t)\right)
\label{eq_interaccion_seguidor_general}
\end{equation}

donde \(a_i(t)\) representa la aceleración del vehículo seguidor, \(v_i(t)\) representa su velocidad actual, \(\Delta v_i(t)\) representa la diferencia de velocidad respecto del líder y \(d_i(t)\) representa la distancia libre disponible. La función \(f(\cdot)\) no se especifica todavía como un modelo particular, sino como una representación conceptual de la dependencia entre el comportamiento del seguidor y las condiciones del vehículo líder. Esta formulación se utiliza como puente hacia los modelos de seguimiento vehicular, donde dicha función puede adoptar expresiones concretas según el modelo seleccionado.

Para una simulación en CARLA, estas variables pueden calcularse a partir del registro de posiciones y velocidades de los vehículos en cada \textit{tick}. El análisis de \(d_i(t)\), \(h_i(t)\), \(\Delta v_i(t)\), \(v_i(t)\) y \(a_i(t)\) permite identificar situaciones de seguimiento cercano, detenciones, aproximaciones rápidas, circulación fluida o posibles condiciones de conflicto. De manera complementaria, estas variables ayudan a explicar por qué una mayor concentración de viajes hacia ciertas zonas puede traducirse en reducción de velocidad, aumento de tiempo de recorrido o mayor frecuencia de interacciones entre vehículos.

\begin{comment}
\subsubsection*{Fuente de las fórmulas}

\begin{itemize}
    \item La ecuación \eqref{eq:posicion_trayectoria} se utiliza como representación funcional de la posición del vehículo en el tiempo. Se adapta a partir de la definición de trayectoria presentada por Immers y Logghe, donde la trayectoria corresponde a la posición del vehículo a través del tiempo.

    \item La ecuación \eqref{eq:trayectoria_discreta} es una adaptación computacional de la trayectoria continua. No aparece de forma literal en la fuente, sino que se expresa como una secuencia de muestras temporales para ajustarla a simulaciones que registran datos por \textit{tick}.

    \item La ecuación \eqref{eq:velocidad_continua} se utiliza como definición estándar de velocidad instantánea. Se reproduce con notación equivalente a la empleada por Immers y Logghe y por la Federal Highway Administration, donde la velocidad individual requiere observación sobre espacio y tiempo.

    \item La ecuación \eqref{eq:velocidad_discreta} es una aproximación por diferencias finitas de la ecuación \eqref{eq:velocidad_continua}. Se adapta al registro discreto de posiciones en simulación computacional.

    \item La ecuación \eqref{eq:aceleracion_continua} se utiliza como definición estándar de aceleración. Se expresa con notación equivalente a la explicación de Immers y Logghe, donde la aceleración corresponde a la segunda derivada de la trayectoria.

    \item La ecuación \eqref{eq:aceleracion_discreta} es una aproximación por diferencias finitas de la aceleración continua. Se adapta al cálculo de aceleración entre dos mediciones consecutivas.

    \item La ecuación \eqref{eq:norma_velocidad} se adapta al formato vectorial de CARLA, cuya API devuelve la velocidad mediante un objeto \texttt{Vector3D}. La fórmula corresponde a la norma euclidiana del vector de velocidad y permite obtener una rapidez escalar.

    \item La ecuación \eqref{eq:space_headway} expresa el \textit{space headway} entre dos vehículos consecutivos con una notación adaptada a un eje longitudinal. Se basa en las definiciones de espaciamiento y \textit{space headway} descritas por Immers y Logghe y por la Federal Highway Administration.

    \item La ecuación \eqref{eq:gap_espacial} separa el espaciamiento total de la distancia libre entre vehículos. Se adapta de la distinción entre longitud física del vehículo y distancia mantenida respecto del vehículo delantero.

    \item La ecuación \eqref{eq:time_headway} se utiliza como definición del \textit{time headway} entre dos vehículos que pasan por una misma sección. Se expresa con notación equivalente a la empleada en teoría de flujo vehicular.

    \item La ecuación \eqref{eq:headway_aproximado} es una simplificación válida bajo velocidad aproximadamente constante y aceleración despreciable durante un intervalo corto. Se adapta a partir de la relación entre espaciamiento, velocidad y separación temporal.

    \item La ecuación \eqref{eq:diferencia_velocidad} se utiliza para representar la velocidad relativa entre vehículo seguidor y vehículo líder. Se expresa como una variable auxiliar para analizar la interacción vehículo-seguidor.

    \item La ecuación \eqref{eq_interaccion_seguidor_general} es una formulación conceptual, no un modelo específico. Se incluye para mostrar que la aceleración del seguidor puede depender de su velocidad, su velocidad relativa y la distancia disponible. Su desarrollo formal se reserva para la sección de modelos microscópicos de comportamiento vehicular.
\end{itemize}
\end{comment}

\section{Modelos microscópicos de comportamiento vehicular}

\subsection{Seguimiento vehicular}

\subsection{Cambio de carril}

\subsection{Intersecciones y semáforos}

\subsection{Reglas de conducción}

\section{Variables macroscópicas del tráfico}

\subsection{Flujo}

\subsection{Densidad}

\subsection{Velocidad media}

\subsection{Ocupación}

\subsection{Relación fundamental $q = ku$}

\section{Diagrama fundamental del tráfico}

\subsection{Flujo libre}

\subsection{Capacidad}

\subsection{Congestión}

\subsection{Modelo de Greenshields}

\section{Capacidad, congestión y ondas de choque}

\subsection{Capacidad vial}

\subsection{Cuellos de botella}

\subsection{Colas}

\subsection{Ondas de choque}

\subsection{Nivel de servicio}

\section{Movilidad urbana y generación de viajes}

\subsection{Concepto de viaje}

\subsection{Demanda de tráfico}

\subsection{Zonas origen-destino}

\subsection{Categorías de destino}

\subsection{Destinos frecuentes y no frecuentes}

\section{Modelos estocásticos de elección de destinos}

\subsection{Variables aleatorias categóricas}

\subsection{Distribuciones ponderadas}

\subsection{Elección jerárquica}

\subsection{Asignación uniforme}

\subsection{Semillas y reproducibilidad}

\subsection{Supuestos de los datos \textit{dummy}}

\section{Redes viales y planificación de rutas}

\subsection{Representación mediante grafos}

\subsection{Costos de recorrido}

\subsection{Dijkstra y A*}

\subsection{Ruta y trayectoria}

\subsection{Asignación de rutas}

\section{Dominio operacional definido}

\subsection{Definición}

\subsection{Condiciones ambientales}

\subsection{Condiciones infraestructurales}

\subsection{Condiciones operacionales}

\subsection{Alcance y limitaciones}

\section{Conducción automatizada y vehículos autónomos}

\subsection{Conceptos generales}

\subsection{Niveles SAE}

\subsection{Percepción, planificación y control}

\subsection{Papel de la simulación}

\section{Simulador CARLA 0.9.15}

\subsection{Propósito}

\subsection{Arquitectura cliente-servidor}

\subsection{Mapas y OpenDRIVE}

\subsection{Actores y \textit{blueprints}}

\subsection{Python API}

\subsection{\textit{Traffic Manager}}

\subsection{\textit{Waypoints}, \textit{spawn points} y rutas}

\subsection{Sensores}

\subsection{Modo síncrono y \textit{no-rendering}}

\subsection{Limitaciones para tráfico}

\section{Datos sintéticos y trazabilidad}

\subsection{Definición de datos sintéticos}

\subsection{Datos de entrada y salida}

\subsection{Reproducibilidad}

\subsection{Calidad de datos}

\section{Análisis del impacto del tráfico}

\subsection{Definición}

\subsection{Escenario base}

\subsection{Escenario propuesto}

\subsection{Variables controladas}

\section{Métricas de análisis de movilidad}

\subsection{Distribución de destinos}

\subsection{Tiempo de viaje}

\subsection{Velocidad media}

\subsection{Tasa de finalización}

\subsection{Incidentes}

\subsection{Invasiones de carril}

\subsection{Concentración de viajes}

\section{Análisis estadístico y visualización}

\subsection{Estadística descriptiva}

\subsection{Comparación de escenarios}

\subsection{Gráficas y visualización}

\endgroup
